\section{Boundary and Phase Effects}

\subsection{Why the admissibility boundary is structurally hard}

Near the admissible boundary $\partial\Omega$, the feasible set is determined by
active inequality constraints $\ThetaSys(\omega)\le 0$ and is therefore
generically non-smooth.
Two structural consequences follow.

\begin{enumerate}
\item \textbf{Directional collapse.}
As constraints become tight, the local cone of admissible directions in the
ambient state space shrinks.
Interventions that would be feasible in the interior lose admissible directions
near $\partial\Omega$, making boundary-crossing actions impossible without
explicit constraint handling or violation.

\item \textbf{Nonlocal propagation.}
Thresholds in OC/ETS are coupled through shared potentials, flows $J$, and
closure cycles $C$.
Enforcing a single active constraint near the boundary typically induces
adjustments in other components of the system, propagating changes nonlocally
through $J$ and $C$ and thereby affecting the continuity measure $k$.
\end{enumerate}

Thus, $\partial\Omega$ is not merely a region of quantitative fragility but a
locus of qualitative structural change in what interventions are admissible.

\subsection{Phase transitions as constraint-induced structural bifurcations}

Even when an intervention is $\ThetaSys$-preserving pointwise, it may alter the
cycles $C$ and flows $J$ required for sustaining a given phase.
Phase transitions in this framework arise not from large perturbations but from
the structural impossibility of maintaining closure conditions under tightened
constraints.

Importantly, such transitions are not removable by arbitrarily reducing the
magnitude of an intervention.
They are induced by the fact that some admissible configurations supporting a
phase cease to exist once specific thresholds become active.
In this sense, phase transitions function as constraint-induced structural
bifurcations rather than dynamical instabilities.

\begin{remark}[Interpretation discipline]
This preprint does not claim that intervention is futile.
It establishes that the combination of boundary-crossing intent and strict
threshold preservation imposes unavoidable structural tradeoffs: any such
intervention must incur either continuity degradation somewhere in $\Omega$ or
a phase change somewhere in its domain.
\end{remark}
